\section{Выбор программно-технических средств реализации}

Для решения задачи разработки веб-приложения автоматизации работы приёмной комиссии БНТУ был выбран стек технологий, включающий фреймворк ASP.NET~Core, систему управления базами данных PostgreSQL, инструмент управления миграциями Liquibase, библиотеку React и компонентную систему Ant Design.
Технологии выбраны с учётом сформулированных в первой главе функциональных и нефункциональных требований, а также с учётом возможностей обеспечения надёжности, поддерживаемости и долгосрочного сопровождения итогового решения силами специалистов, доступных в среде учреждения высшего образования.

\subsection{Архитектурное решение}

Прежде чем перейти к выбору конкретных программных средств, целесообразно зафиксировать общее архитектурное решение, которому будут подчинены остальные выборы.

Приложение реализовано в форме клиент-серверной системы с разделением на две независимо разрабатываемые и развёртываемые части: одностраничное веб-приложение (Single Page Application), исполняющееся в браузере пользователя, и серверную часть, обрабатывающую запросы клиента и взаимодействующую с базой данных.
Подобное разделение является устоявшимся подходом в современной веб-разработке и обеспечивает независимость технологического стека клиента и сервера, упрощает разработку, развёртывание и масштабирование частей системы по отдельности.

Взаимодействие клиентской и серверной частей построено на принципах архитектурного стиля REST поверх протокола HTTP с обменом данными в формате JSON.
Этот формат хорошо поддерживается всеми современными платформами разработки, удобен как для машинной обработки, так и для отладки, и не требует подключения дополнительных библиотек сериализации.

В качестве модели данных выбрана реляционная модель.
Этот выбор естественно соответствует природе предметной области: правила приёма представлены набором взаимосвязанных сущностей (специальности, конкурсные списки, основания зачисления, показатели оценивания, группы показателей, заявления абитуриентов), целостность связей между которыми должна обеспечиваться на уровне системы хранения.
Реляционные системы управления базами данных предоставляют для этой цели зрелый аппарат: внешние ключи с каскадными операциями, проверочные ограничения, уникальные составные ключи, транзакционные гарантии ACID при изменении нескольких связанных таблиц.

Для аутентификации пользователей выбран подход на основе токенов JSON Web Token, подписываемых на серверной стороне.
В отличие от подхода с серверными сессиями, такой механизм не требует хранения состояния сессий на сервере, что упрощает горизонтальное масштабирование серверной части и устраняет необходимость в общем хранилище сессий.
Токен передаётся клиентом в заголовке \code{Authorization} каждого запроса и проверяется серверной частью без обращения к базе данных.

\subsection{Серверная часть: ASP.NET Core}

ASP.NET~Core~--~кросс-платформенный фреймворк для разработки веб-приложений и сервисов на платформе~.NET, разрабатываемый компанией Microsoft и распространяемый под открытой лицензией MIT.
Фреймворк выбран для серверной части по следующим причинам.

Кросс-платформенность.
Поддержка операционных систем семейств Windows и Linux обеспечивает гибкость при развёртывании в инфраструктуре университета и снижает требования к лицензированию серверного программного обеспечения.

Встроенные механизмы безопасности.
ASP.NET~Core предоставляет готовые компоненты для аутентификации с использованием JSON Web Tokens, авторизации на основе атрибутов и политик, валидации входных данных через атрибуты Data Annotations, что снижает класс ошибок, связанных с самостоятельной реализацией типовых механизмов безопасности.

Поддержка REST API и сериализации.
Контроллеры с маршрутизацией на основе атрибутов позволяют декларативно описывать REST-эндпоинты, автоматически выполнять сериализацию объектов в формат JSON и обратное преобразование, а также предоставлять интерактивную документацию через Swagger в режиме разработки.

Зрелость экосистемы.
Платформа~.NET имеет развитую экосистему библиотек, документации и инструментов разработки, в том числе официально поддерживаемый драйвер Npgsql для PostgreSQL, библиотеку ClosedXML для работы с форматом Excel и библиотеку BCrypt.Net-Next для безопасного хэширования паролей.

В качестве серверной платформы для веб-приложений на сегодняшний день наиболее распространены ASP.NET~Core (язык программирования C\#), Spring~Boot (Java), Node.js с фреймворками Express или NestJS (JavaScript и TypeScript) и Django (Python).
Каждое из перечисленных решений обеспечивает требуемый функционал: маршрутизацию HTTP-запросов, сериализацию данных, аутентификацию и работу с базой данных.

Выбор ASP.NET~Core среди этих альтернатив обусловлен соображениями долгосрочной поддерживаемости разработанной системы.
Язык C\# и платформа~.NET широко представлены в учебных программах Белорусского национального технического университета, в том числе на факультете информационных технологий и робототехники, и используются при подготовке специалистов по информационным технологиям.
Это означает, что задачу сопровождения и развития системы в дальнейшем может выполнять как штатный сотрудник из числа преподавателей профильных кафедр, так и студент старших курсов в рамках производственной практики или дипломной работы, без необходимости привлечения внешних специалистов с редкой для университетской среды квалификацией.

\subsection{Прямые SQL-запросы через Npgsql}

В рамках данного проекта принято решение отказаться от использования объектно-реляционного преобразователя (ORM) и работать с базой данных напрямую через библиотеку Npgsql~--~официальный драйвер PostgreSQL для платформы~.NET, распространяемый под лицензией PostgreSQL License.
Решение обусловлено следующими соображениями.

Несоответствие объектно-реляционной парадигмы природе доменной модели проекта.
Доменная модель (раскрывается подробно в третьей главе) представлена настраиваемыми сущностями: показатели оценивания абитуриентов, группы показателей, основания зачисления и связи между ними хранятся в базе данных как данные, а не как фиксированный набор классов в исходном коде.
В таком контексте объектно-реляционный преобразователь не даёт характерных для него преимуществ: отображение <<таблица~$\leftrightarrow$~класс>> неприменимо, поскольку набор <<классов>> предметной области определяется не кодом приложения, а содержимым справочников и меняется во времени без перекомпиляции.

Контроль над генерируемыми запросами.
Алгоритм пересчёта результатов отбора требует загрузки сводного снимка данных одним запросом с агрегированными результатами, что неэффективно реализуется средствами объектно-реляционных преобразователей и приводит к избыточным запросам типа N+1.

Применение специфичных возможностей PostgreSQL.
Использование оператора \code{ON CONFLICT}, операции \code{COPY BINARY} для массовой вставки результатов отбора и явных блокировок \code{pg\_advisory\_xact\_lock} требует написания SQL-запросов напрямую и средствами объектно-реляционных преобразователей выражается ограниченно.

Простота отладки.
Прямой SQL-код легко переносится в любой инструмент администрирования базы данных для воспроизведения проблемы и анализа плана выполнения запроса, что упрощает диагностику в эксплуатации.

\subsection{Система управления базами данных PostgreSQL}

PostgreSQL~--~свободно распространяемая объектно-реляционная система управления базами данных, разрабатываемая международным сообществом и распространяемая под лицензией PostgreSQL License, совместимой с BSD/MIT.
PostgreSQL выбрана в качестве системы управления базами данных по следующим причинам.

Свободное распространение.
Лицензия PostgreSQL License не накладывает финансовых обязательств на учреждение высшего образования и не требует раскрытия исходного кода приложений, использующих эту СУБД, что упрощает развёртывание системы.

Поддержка сложных типов данных.
PostgreSQL поддерживает массивы, JSON, перечислимые типы (\code{ENUM}), типы для работы с диапазонами значений.
В проекте перечислимые типы используются, в частности, для хранения типа показателя оценивания (значения \code{higher\_is\_better} и \code{lower\_is\_better}) и статусов запросов на удаление абитуриентов.

Развитая система ограничений целостности.
Поддержка проверочных ограничений, уникальных составных ключей и внешних ключей с каскадными операциями обеспечивает целостность данных на уровне СУБД и не полагается исключительно на корректность приложения.
Это особенно важно для проекта с настраиваемой доменной моделью: целостность ссылок между взаимосвязанными настраиваемыми сущностями (показателями, группами, основаниями зачисления) контролируется на уровне СУБД и сохраняется даже при ошибках в коде приложения.

Производительная работа с большими объёмами данных.
Возможность создания частичных индексов, выражение-индексов и расширенный планировщик запросов позволяют эффективно обрабатывать запросы постраничного получения списков из таблиц с десятками тысяч записей.

\subsection{Управление миграциями: Liquibase}

Для управления версиями схемы базы данных выбран инструмент Liquibase в свободной редакции Community, распространяемой под лицензией Apache~2.0.
Liquibase позволяет описывать изменения схемы в виде набора файлов формата YAML и SQL, отслеживать применённые изменения через служебную таблицу в базе данных и автоматически применять недостающие изменения при обновлении приложения.

По сравнению с инструментом миграций Entity Framework Core, Liquibase обеспечивает декларативный подход с возможностью просмотра истории изменений в виде отдельных файлов, не требует наличия объектно-реляционного преобразователя на стороне приложения (что согласуется с описанным выше отказом от ORM) и не привязан к конкретной платформе разработки, благодаря чему может быть переиспользован при возможной миграции серверной части на иной технологический стек.

\subsection{Клиентская часть: React}

React~--~библиотека для построения пользовательских интерфейсов на JavaScript, разрабатываемая компанией Meta и распространяемая под лицензией MIT.
Библиотека выбрана для разработки клиентской части по следующим причинам.

Компонентная архитектура.
Декомпозиция интерфейса на повторно используемые компоненты с собственным состоянием и поведением упрощает разработку и сопровождение сложных пользовательских интерфейсов, в частности страниц ведения справочников настраиваемой доменной модели.

Декларативный подход.
Состояние компонента описывается через набор переменных, а представление автоматически перерисовывается при их изменении, что снижает количество ошибок, связанных с рассинхронизацией состояния и интерфейса.

Богатая экосистема.
Для библиотеки React существует широкий выбор готовых решений: React~Router для маршрутизации, Axios для HTTP-запросов, react-i18next для интернационализации, что ускоряет разработку и снижает потребность в собственной реализации типовых задач.

Использование инструмента сборки Vite (лицензия MIT) обеспечивает быструю разработку благодаря технологии Hot Module Replacement и оптимизированную сборку для производственной среды.

Среди современных библиотек и фреймворков построения одностраничных веб-приложений основными альтернативами являются React, Angular и Vue.js.
Все три решения зрелые, активно развиваются и предоставляют необходимый функционал: компонентную модель, маршрутизацию, управление состоянием.

React выбран как наиболее распространённое из перечисленных решений: по данным регулярных опросов разработчиков (Stack Overflow Developer Survey, State of JavaScript) и статистике установок пакетов в реестре npm, React занимает первое место по доле использования среди библиотек и фреймворков клиентской веб-разработки на протяжении последнего десятилетия.
Высокая распространённость React означает обширное сообщество, обилие учебных материалов в открытом доступе (в том числе на русском языке), большое количество готовых решений для типовых задач и низкий порог входа для нового разработчика, который будет привлечён к сопровождению системы.
Эти соображения симметричны соображениям выбора серверной платформы и подчинены общей задаче обеспечения непрерывности сопровождения разработанной системы силами специалистов общедоступной квалификации.

\subsection{Компонентная система Ant Design}

Ant Design~--~библиотека компонентов пользовательского интерфейса корпоративного класса, разрабатываемая компанией Ant Group и распространяемая под лицензией MIT.
Библиотека предоставляет более шестидесяти готовых компонентов: таблицы с серверной пагинацией, формы со встроенной валидацией, модальные окна, навигационные меню и другие.
К преимуществам Ant Design относятся: единая визуальная стилистика всех компонентов, поддержка тёмной и светлой темы, локализация на множество языков (в том числе русский и английский), развитая система настройки внешнего вида через токены дизайна.

В проекте Ant Design выступает основой для интерфейса администратора, через который и происходит конфигурирование доменной модели правил приёма: компонент \code{Form} со встроенной валидацией используется для редактирования справочников, компонент \code{Table} с серверной пагинацией~--~для отображения списков записей, компонент \code{Modal} с подтверждением~--~для безопасного выполнения деструктивных действий.

\subsection{Вспомогательные библиотеки}

Помимо перечисленных, в проекте используются следующие библиотеки.

Axios~--~библиотека для выполнения HTTP-запросов с поддержкой интерцепторов запроса и ответа (лицензия MIT).
Применяется на клиенте для централизованного добавления заголовков авторизации и языка интерфейса ко всем исходящим запросам, что устраняет необходимость дублировать соответствующий код в каждом месте обращения к серверной части.

ClosedXML~--~библиотека для платформы~.NET для работы с файлами формата Office~Open~XML Spreadsheet (\code{xlsx}), распространяемая под лицензией MIT.
Используется на сервере для формирования файлов с результатами отбора, передаваемых сотрудникам приёмной комиссии для дальнейшей обработки и, при необходимости, для последующего внесения данных во внешние системы.

BCrypt.Net-Next~--~реализация алгоритма хэширования паролей bcrypt для платформы~.NET (лицензия MIT).
Применяется при создании учётных записей пользователей и проверке пароля при входе: пароль преобразуется в хэш с автоматически сгенерированной солью, что исключает хранение паролей в открытом виде и обеспечивает устойчивость к атакам подбора по предвычисленным таблицам.

react-i18next~--~библиотека интернационализации для React, обеспечивающая загрузку строк перевода из файлов формата JSON, автоматическое определение языка пользователя и переключение языка без перезагрузки страницы (лицензия MIT).

Microsoft.Extensions.Localization~--~встроенный в платформу~.NET механизм локализации, обеспечивающий загрузку строк перевода из файлов ресурсов \code{.resx} и их использование как через атрибуты валидации, так и напрямую через инжектируемый интерфейс \code{IStringLocalizer}.

\subsection{Лицензионная характеристика выбранного стека}

Выбранный стек технологий обеспечивает реализацию сформулированных в первой главе функциональных и нефункциональных требований и обладает следующей общей особенностью.

Все использованные в проекте библиотеки, фреймворки и инструменты распространяются под пермиссивными свободными лицензиями: MIT, Apache~2.0 и PostgreSQL License.
Эти лицензии не накладывают ограничений на коммерческое использование разработанной системы, не требуют раскрытия её исходного кода при использовании соответствующих библиотек и совместимы между собой при совместном применении.
Это означает отсутствие у учреждения высшего образования каких-либо финансовых или организационных обязательств перед правообладателями использованных программных компонентов и снимает юридические препятствия к развёртыванию системы в инфраструктуре университета.
